\documentclass[10pt,a4paper]{amsart}
\usepackage[margin=19mm]{geometry}
\usepackage{amsmath,amssymb,mathtools}
\usepackage[scr=rsfs,cal=euler]{mathalfa}
\usepackage{xfrac}
\usepackage[unicode,hidelinks,bookmarks=false]{hyperref}
\newcommand{\ie}{i.e.\ }
\newcommand{\cf}{cf.\ }
\newcommand{\resp}{respectively\ }
\newcommand{\Subsection}{\subsection}
\providecommand{\nobreakdash}{\nobreak\hbox{-}}
\setlength{\emergencystretch}{2em}
\multlinegap0pt
\usepackage{mathtools}
\usepackage{amssymb}
\usepackage{float}
\usepackage{tikz}
\newtheorem{lemma}{Lemma}
\newtheorem{corollary}{Corollary}
\newcommand{\C}{\mathcal{C}}
\def\N{\mathbb{N}}
\title{On the Tur\'anability and tileability of oriented graphs}
\author{Igor Araujo, Zimu Xiang}
\date{Source: arXiv:2507.13267v2 (CC BY 4.0)}
\begin{document}
\maketitle

\noindent\emph{Source and licence.} On the Tur\'anability and tileability of oriented graphs. Version arXiv:2507.13267v2 (CC BY 4.0), distributed by arXiv under the Creative Commons Attribution 4.0 International licence, http://creativecommons.org/licenses/by/4.0/, as recorded in the arXiv licence metadata for this exact version and shown on the abstract page https://arxiv.org/abs/2507.13267v2. Full licence notice: \texttt{licenses/arxiv-2507.13267-CC-BY-4.0.txt}.\par\smallskip

\noindent\textit{Editorial scope.} Counted unit: the article's lemma lem:extremal (source0.tex lines 725--728). The article's complete original proof of it is delivered with it (source0.tex lines 952--1007, one \texttt{proof} environment of 56 lines, which the article places away from the statement). No mathematics is written, repaired or re-derived: the statement and the proof are the article's own lines, in the article's own notation, and no retained line is substituted or re-flowed.  Retained uncounted as the source-local context the proof needs: lines 753--755; lines 888--898; lines 930--932.  Nothing was omitted from any retained span, so the package carries no omission record.  No retained line contains a citation, so the wrapper carries no bibliography and no bibliography entry.  No retained line contains a figure environment, an image-inclusion command or drawing code, so no image is delivered and the package declares no asset dependency.  Editorial formatting only, in addition: an AMS \texttt{amsart} preamble that restates the article's own declarations and macros used by the retained lines, namely the environments \texttt{lemma}, \texttt{corollary} on separate counters, together with the standard packages those lines need.  The retained environments print in this document as Lemma 1, Corollary 1 in order of appearance, whereas the article prints the same items with the numbering of its own full document.  The complete native source of the article as distributed in the arXiv e-print for this exact version is bundled unchanged as \texttt{sources/181759/source0.tex}.
\begin{lemma}\label{lem:minimum_degree}
    Suppose that $c>0$ and $G$ is an oriented graph on $n$ vertices with $\delta^0(G)\ge (1/2-c)n$, then each vertex of $G$ belongs to at least $(1/32-4c)n^3$ copies of $D$.
\end{lemma}

\begin{corollary}\label{cor:0mod4}
    There exist $c', \xi >0$ and $n_0\in \N$ so that the following holds. Suppose $n\ge n_0$ is a multiple of $4$, and $G$ is a near-tournament on $n$ vertices such that $\delta^0(G)\ge (1/2-c')n$. If there is a partition $\{V_1,V_2,V_3\}$ of $V(G)$ such that 
    \begin{equation} \label{eq:div}
        |V_1|\equiv |V_2|\equiv |V_3|\equiv 0\pmod 4
    \end{equation}
    and, for every $i\in [3]$ and $v\in V_i$, 
    \begin{equation} \label{eq:v}
        d^+(v,V_{i+1}), \; d^-(v,V_{i-1})\ge (1/3-\xi)n ,
    \end{equation}
    then $G$ has a perfect $D$-tiling.
\end{corollary}

\begin{lemma}\label{lem:ith}
    Let $\alpha, \beta, c>0$ be such that $\alpha+2\beta+2c < 1/3$. Let $G$ be an oriented graph with $\delta^0(G)\ge (1/2-c)n$. Suppose that there are disjoint $V_1,V_2,V_3\subseteq V(G)$ satisfying that for every $i\in [3]$, we have $|V_i|\ge (1/3-\alpha)n$ and $V_i\subseteq N^-_\beta(V_{i+1})$ and $V_i\subseteq N^+_\beta(V_{i-1})$. Then, for every $u\in V_1\cup V_2\cup V_3$ and $i \in [3]$, there is a copy of $D$ in $G$ of $i$-th type that contains the vertex $u$.
\end{lemma}

\begin{lemma}[Extremal case] \label{lem:extremal}
    There exist $c, \gamma >0$ and $n_0\in \N$ so that the following holds. Suppose that $n\ge n_0$ is divisible by $4$ and $G$ is a near-tournament on $n$ vertices. 
    If $\delta^0(G)\ge (1/2-c)n$ and $G$ is $\gamma$-extremal, then $G$ has a perfect $D$-tiling.
\end{lemma}

\begin{proof}[Proof of Lemma~\ref{lem:extremal}]
    In what follows, we assume that $c, \gamma >0$ are sufficiently small so that there exists $\beta > \max\{ 4c, \sqrt{60\gamma} \}$ satisfying   
    \begin{equation}\label{eq:beta<c}
        \beta < \frac{1}{32}-4c
    \end{equation}
    \begin{equation}\label{eq:bgc<1/3}
        9\beta + \gamma + 2c < \frac{1}{3}
    \end{equation}
    and so that Corollary~\ref{cor:0mod4} holds with parameters $c' > c+3\beta$ and $\xi > 2\beta +\gamma $. 
    
    We fix a partition $\{V_1,V_2,V_3\}$ of $V(G)$ such that, for every $i\in[3]$, $(1/3 - \gamma)n \le |V_i| \le (1/3 + \gamma)n$ and $e^-(V_i,V_{i+1})\le \gamma n^2$.
    Let $U_i=N^-_\beta(V_{i+1})\cap N^+_\beta(V_{i-1}) \cap V_i$ for $i\in [3]$, and $U_0=V\setminus(U_1\cup U_2\cup U_3)$.
    For every $i\in[3]$ and $v\in V_i\setminus U_i$, by definition of $U_i$, at least one of the following holds:
    \begin{align*}
        d^-(v,V_{i+1})\ge |V_{i+1}|-(|V_{i+1}|-\beta n)-(n-2\delta^0(G))=\beta n-2cn > \beta n/2,\\
        d^+(v,V_{i-1})\ge |V_{i-1}|-(|V_{i-1}|-\beta n)-(n-2\delta^0(G))=\beta n-2cn > \beta n/2.
    \end{align*}
    Thus we have that
    \begin{equation}
        |U_0|=\sum_{i=1}^3|V_i\setminus U_i|\le \frac{\sum_{i=1}^3 2e^-(V_i,V_{i+1})}{\beta n/2}\le \frac{12\gamma n^2}{\beta n} < \frac{\beta n}{5}.\label{equ:u0}
    \end{equation}

    By~\eqref{eq:beta<c},~\eqref{equ:u0} and Lemma~\ref{lem:minimum_degree}, we can greedily find a collection $\C_0$ of vertex-disjoint copies of $D$ that covers all of $U_0$.
    For $i\in[3]$, let $W_i=U_i\setminus V(\C_0)$ and $W=W_1\cup W_2\cup W_3$.
    By~\eqref{equ:u0}, since $|V(\C_0)|\le 4|U_0|$, we have that $|V_i\setminus W_i|\le |V(\C_0)| + |U_0| < \beta n$. This implies that
    \begin{equation} \label{eq:|Wi|}
        |W_i| > |V_i| - \beta n \ge \left(\frac{1}{3} - \gamma - \beta \right) n, \text{ and}
    \end{equation}
    \begin{equation} \label{eq:deg_Wi}
        \delta^0(G[W]) > \left( \frac{1}{2} -c \right)n - 3\beta n \ge 
        \left( \frac{1}{2} -c - 3\beta \right) |W|.
    \end{equation}

    For $v \in W_i$, we have that $v \in N_\beta^-(V_{i+1})$, which implies $d^+(v, V_{i+1}) \ge |V_{i+1}|-\beta n$ and then $d^+(v, W_{i+1}) \ge |W_{i+1}|-\beta n$. Similarly, $d^-(v, W_{i-1}) \ge |W_{i-1}|-\beta n$. That is, 
    \begin{equation} \label{eq:Wi_subset}
        W_i \subseteq N_\beta^-(W_{i+1}) \cap N_\beta^+(W_{i-1}).
    \end{equation}
    
    Let the triple $(a_1,a_2,a_3)\in \{0,1,2,3\}^3$ be such that $|W_i| \equiv a_i \pmod 4$ for every $i \in [3]$. Note that $a_1+a_2+a_3 \equiv n \equiv 0 \pmod 4$. For every $i \in [3]$, set $b_i \equiv a_i - \frac{a_1+a_2+a_3}{4} \pmod 4$ so that $b_i \in \{0,1,2,3\}$.
    
    By~\eqref{eq:bgc<1/3},~\eqref{eq:|Wi|},~\eqref{eq:deg_Wi},~\eqref{eq:Wi_subset}, and Lemma~\ref{lem:ith} applied multiple times, there exists a collection $\C_1$ of disjoint copies of $D$ such that $\C_1$ contains exactly $b_i$ many copies of $D$ of the $i$-th type for every $i\in [3]$.
    It follows
    \begin{itemize}
        \item For $W_i'=W_i\setminus V(\C_1)$, we have $|W_i'| \equiv a_i - 2b_i - b_{i+1} -b_{i+2}\equiv 0\pmod 4$,
        \item $|V(\C_1)| = 4 (b_1+b_2+b_3) \le 36$.
    \end{itemize}

    Let $W'=W_1'\cup W_2'\cup W_3'$. We note that, for every $i\in [3]$ and $v \in W_i'$,
    \begin{equation}
        d^+(v, W_{i+1}') \ge d^+(v, W_{i+1}) - 12 \ge   
        \left( \frac{1}{3} - \gamma - 2\beta \right) n \ge
        \left( \frac{1}{3} - \xi \right) |W'|,
    \end{equation}
    and, similarly, $d^-(v, W_{i-1}')\ge \left( \frac{1}{3} - \xi \right) |W'|$.
    Then, by Corollary~\ref{cor:0mod4}, there is some perfect $D$-tiling $\C_2$ of $G[W']$, and $\C_0 \cup\C_1\cup \C_2$ is a perfect $D$-tiling of $G$.
\end{proof}

\end{document}
